% !TEX encoding = UTF-8 Unicode
\documentclass[12pt,a4paper]{article}
\usepackage[utf8]{inputenc}
\usepackage[T1]{fontenc}
\usepackage[ngerman,latin,french]{babel}
\usepackage{amsmath,amssymb,amsthm}
\usepackage{geometry}
\usepackage{titlesec}
\usepackage{fancyhdr}
\usepackage{microtype}
\usepackage{longtable}
\usepackage{array}

\geometry{
  left=2.5cm,
  right=2.5cm,
  top=2.5cm,
  bottom=2.5cm
}

\pagestyle{fancy}
\fancyhf{}
\fancyhead[C]{\small\nouppercase{\leftmark}}
\fancyfoot[C]{\thepage}
\renewcommand{\headrulewidth}{0.4pt}

\titleformat{\section}{\large\bfseries}{\thesection.}{0.5em}{}
\titleformat{\subsection}{\normalsize\bfseries}{\thesubsection.}{0.5em}{}

\newcommand{\dd}{\mathop{}\!\mathrm{d}}
\newcommand{\D}{\Delta}

\title{Carl Friedrich Gau\ss{} Werke --- Band 7, Teil 10}
\author{Carl Friedrich Gau\ss{}}
\date{}

\begin{document}

\maketitle
\thispagestyle{empty}

\tableofcontents
\newpage

\selectlanguage{french}
\section*{Exposition d'une nouvelle m\'{e}thode de calculer les perturbations plan\'{e}taires avec l'application au calcul num\'{e}rique des perturbations du mouvement de Pallas.}
\addcontentsline{toc}{section}{Exposition d'une nouvelle m\'{e}thode}

\subsection*{\S\,3.}
\addcontentsline{toc}{subsection}{\S\,3}

Ainsi en faisant abstraction des perturbations nous aurons
\begin{align}
\frac{\dd^2 x}{\dd t^2} &= -\frac{mx}{r^3}, \tag{1} \\
\frac{\dd^2 y}{\dd t^2} &= -\frac{my}{r^3}, \tag{2} \\
\frac{\dd^2 z}{\dd t^2} &= -\frac{mz}{r^3}, \tag{3}
\end{align}
o\`{u} $r$ est le rayon vecteur de la plan\`{e}te au soleil.

Il s'ensuit, que faisant
\begin{equation}
\alpha = \frac{x\,\dd y - y\,\dd x}{\dd t}, \quad
\beta = \frac{y\,\dd z - z\,\dd y}{\dd t}, \quad
\gamma = \frac{z\,\dd x - x\,\dd z}{\dd t}, \tag{4}
\end{equation}
on aura
\begin{equation}
\frac{\dd \alpha}{\dd t} = 0, \quad \frac{\dd \beta}{\dd t} = 0, \quad \frac{\dd \gamma}{\dd t} = 0, \tag{5}
\end{equation}
d'o\`{u} $A$, $B$, $C$ seront des quantit\'{e}s constantes, qui par cons\'{e}quent d\'{e}pendront de l'\'{e}tat primitif.

On tire de ces \'{e}quations
\begin{equation}
Ax + By + Cz = 0. \tag{6}
\end{equation}

L'orbite sera donc dans un plan passant par le centre du soleil. Nous avons de plus
\begin{align*}
r^2 &= x^2 + y^2 + z^2, \\
r\,\dd r &= x\,\dd x + y\,\dd y + z\,\dd z, \\
r\,\dd^2 r + \dd r^2 &= \dd x^2 + \dd y^2 + \dd z^2.
\end{align*}

Or les \'{e}quations (1), (2), (3) donnent
\begin{multline*}
(AA + BB + CC)\,\dd t^2 = (x^2 + y^2 + z^2)(\dd x^2 + \dd y^2 + \dd z^2) \\
- (x\,\dd x + y\,\dd y + z\,\dd z)^2 = r^2(\dd x^2 + \dd y^2 + \dd z^2) - r^2\,\dd r^2 = r^3\,\dd^2 r + mr\,\dd t^2.
\end{multline*}

Ainsi en faisant
\[
\frac{AA + BB + CC}{m} = p, \quad r - p = u,
\]
on a
\begin{equation}
r\,\dd^2 r + \dd r^2 = \frac{p}{r}\,\dd t^2, \tag{7}
\end{equation}
\'{e}quation analogue aux trois premi\`{e}res. On voit donc qu'en faisant
\begin{align}
\frac{x\,\dd u - u\,\dd x}{\dd t} &= F, &
\frac{y\,\dd u - u\,\dd y}{\dd t} &= G, &
\frac{z\,\dd u - u\,\dd z}{\dd t} &= H, \tag{8}
\end{align}
$F$, $G$, $H$ seront aussi des quantit\'{e}s constantes.

La somme des \'{e}quations (8) multipli\'{e}es par $x$, $y$, $z$ nous donne
\begin{equation}
Fx + Gy + Hz = \frac{m}{r}(r-p). \tag{9}
\end{equation}

Multipliant celle-ci par $F$, (8) par $H$, (8) par $-G$, et ajoutant, il vient
\[
(FF + GG + HH)x = (BH - CG)(r-p) + \frac{Fr\,\dd r}{\dd t},
\]
o\`{u} \`{a} cause de $FF + GG + HH = (1-e^2)mp = mpee$:

\begin{equation}
x = \frac{a(BH-CG)}{mpe} \cos E + \frac{F \sin E}{\sqrt{am}}. \tag{10}
\end{equation}

On aura de m\^{e}me
\begin{align}
y &= \frac{a(CF-AH)}{mpe} \cos E + \frac{G \sin E}{\sqrt{am}}, \tag{11} \\
z &= \frac{a(AG-BF)}{mpe} \cos E + \frac{H \sin E}{\sqrt{am}}. \tag{12}
\end{align}

Or on a
\[
r - p = aee - ae\cos E, \quad r\,\dd r = ae\sin E \cdot \dd E = \sqrt{am} \cdot ae\sin E \cdot \dd t;
\]
ainsi nous aurons
\begin{align}
x &= \frac{a(BH-CG)}{mpe} \cos E + \frac{F \sin E}{\sqrt{am}}, \tag{13} \\
y &= \frac{a(CF-AH)}{mpe} \cos E + \frac{G \sin E}{\sqrt{am}}, \tag{14} \\
z &= \frac{a(AG-BF)}{mpe} \cos E + \frac{H \sin E}{\sqrt{am}}. \tag{15}
\end{align}

Les quatre \'{e}quations (9), (13), (14), (15) contiennent la solution compl\`{e}te des \'{e}quations (1), (2), (3).

\subsection*{\S\,4.}
\addcontentsline{toc}{subsection}{\S\,4}

Il nous reste encore de d\'{e}velopper la liaison entre nos constantes et les \'{e}l\'{e}mens dont les astronomes ont coutume de faire usage. Pour cet effet, rapportons \`{a} une surface sph\'{e}rique d'un rayon arbitraire et dont le centre co\"{i}ncide avec celui du soleil, les neuf directions suivantes: les trois axes des coordonn\'{e}es; la droite men\'{e}e du soleil vers la plan\`{e}te; une droite perpendiculaire \`{a} celle-ci dans le plan de l'orbite; la droite perpendiculaire \`{a} ce plan; la ligne des apsides men\'{e}e vers le p\'{e}rih\'{e}lie; la perpendiculaire \`{a} celle-ci dans le plan de l'orbite; l'intersection du plan des $x$, $y$ avec le plan de l'orbite. Nous d\'{e}noterons les points de la surface sph\'{e}rique auxquels r\'{e}pondent ces directions, par $X$, $Y$, $Z$, $L$, $A$, $M$, $P$, $\Pi$, $N$. Ainsi $X$, $Y$, $Z$ seront les p\^{o}les des grands cercles qui repr\'{e}sentent les trois plans fondamentaux, et nous supposerons ces p\^{o}les pris du c\^{o}t\'{e} o\`{u} les coordonn\'{e}es sont cens\'{e}es positives; $L$ sera le lieu h\'{e}liocentrique de la plan\`{e}te et $A$ le point du grand cercle de l'orbite avanc\'{e} davantage que $L$ de $90^\circ$; $P$ le p\'{e}rih\'{e}lie; $\Pi$ ce que devient $A$ lorsque $L$ devient $P$; $N$ sera le n\oe ud ascendant de l'orbite sur le grand cercle dont $Z$ est le p\^{o}le: enfin $M$ sera l'un des p\^{o}les de l'orbite, et pour ne laisser aucune ambigu\"{\i}t\'{e} nous prendrons pour $M$ celui qui est du c\^{o}t\'{e} o\`{u} les coordonn\'{e}es $z$ sont positives.
\newpage
\subsection*{\S\,5.}
\addcontentsline{toc}{subsection}{\S\,5}

Nous d\'{e}signerons l'arc $XN$ compt\'{e} de $X$ vers $Y$, par $\Omega$, et la somme de cet arc avec l'arc $NL$ compt\'{e} de $N$ dans le sens du mouvement par $v$; cette somme est ce que les astronomes nomment la longitude dans l'orbite. Soit $\omega$ ce que devient $v$ lorsque $L$ devient $P$, ou la longitude du p\'{e}rih\'{e}lie; $v - \omega$ sera donc l'arc $PL$ ou l'anomalie vraie. L'inclinaison du plan dans lequel se meut la plan\`{e}te sur la branche du grand cercle $XY$ qui va de $N$ dans le sens direct sera \'{e}gale \`{a} $ZM$; nous la d\'{e}signerons par $i$. Cela pos\'{e} les 15 formules pr\'{e}c\'{e}dentes peuvent \^{e}tre repr\'{e}sent\'{e}es de la mani\`{e}re suivante:
\begin{align}
x &= r(\cos\Omega \cdot \cos(v - \Omega) - \sin\Omega \cdot \sin(v - \Omega) \cdot \cos i), \tag{16} \\
y &= r(\sin\Omega \cdot \cos(v - \Omega) + \cos\Omega \cdot \sin(v - \Omega) \cdot \cos i), \tag{17} \\
z &= r\sin(v - \Omega)\sin i, \tag{18} \\
A &= \sqrt{mp} \cdot \sin\Omega \cdot \sin i, \tag{19} \\
B &= -\sqrt{mp} \cdot \cos\Omega \cdot \sin i, \tag{20} \\
C &= \sqrt{mp} \cdot \cos i, \tag{21} \\
Bz - Cy &= r\sqrt{mp} \cdot (-\cos\Omega \cdot \sin(v - \Omega) - \sin\Omega \cdot \cos(v - \Omega) \cdot \cos i), \tag{22} \\
Cx - Az &= r\sqrt{mp} \cdot (-\sin\Omega \cdot \sin(v - \Omega) + \cos\Omega \cdot \cos(v - \Omega) \cdot \cos i), \tag{23} \\
Ay - Bx &= r\sqrt{mp} \cdot \cos(v - \Omega) \cdot \sin i, \tag{24} \\
CG - BH &= mpe(\cos\Omega \cdot \cos(\omega - \Omega) - \sin\Omega \cdot \sin(\omega - \Omega) \cdot \cos i), \tag{25} \\
AH - CF &= mpe(\sin\Omega \cdot \cos(\omega - \Omega) + \cos\Omega \cdot \sin(\omega - \Omega) \cdot \cos i), \tag{26} \\
BF - AG &= mpe\sin(\omega - \Omega)\sin i, \tag{27} \\
F &= e\sqrt{mp} \cdot (-\cos\Omega \cdot \sin(\omega - \Omega) - \sin\Omega \cdot \cos(\omega - \Omega) \cdot \cos i), \tag{28} \\
G &= e\sqrt{mp} \cdot (-\sin\Omega \cdot \sin(\omega - \Omega) + \cos\Omega \cdot \cos(\omega - \Omega) \cdot \cos i), \tag{29} \\
H &= e\sqrt{mp} \cdot \cos(\omega - \Omega) \cdot \sin i. \tag{30}
\end{align}

Multipliant les \'{e}quations (10), (11), (12) par $A$, $B$, $C$ et ajoutant, on trouve
\[
x(BH - CG) + y(CF - AH) + z(AG - BF) = mp(r - p).
\]

Or le premier membre de cette \'{e}quation devient, moyennant les \'{e}quations (16)--(18), (25)--(27)
\[
= -mper\cos PL = -mper\cos(v - \omega);
\]
ainsi on a
\begin{equation}
r - p = -er\cos(v - \omega), \quad \text{donc} \quad \frac{p}{r} = 1 + e\cos(v - \omega). \tag{31}
\end{equation}

De l\`{a} on tire
\begin{equation}
\frac{1}{r} = \frac{1 + e\cos(v - \omega)}{p}, \quad \text{et} \quad r = \frac{p}{1 + e\cos(v - \omega)}. \tag{32}
\end{equation}

\newpage
\subsection*{\S\,6.}
\addcontentsline{toc}{subsection}{\S\,6}

D\'{e}signons par $R$ le rayon vecteur de la plan\`{e}te troublee dans l'orbite osculatrice, et par $\varphi$ l'anomalie vraie compt\'{e}e du p\'{e}rih\'{e}lie, de sorte que
\[
R = \frac{p}{1 + e\cos\varphi}, \quad \text{ou} \quad \frac{p}{R} = 1 + e\cos\varphi.
\]

On a dans l'orbite elliptique
\begin{equation}
\frac{\dd^2 R}{\dd t^2} - R\left(\frac{\dd\varphi}{\dd t}\right)^2 = -\frac{m}{R^2}, \quad R\frac{\dd^2\varphi}{\dd t^2} + 2\frac{\dd R}{\dd t}\frac{\dd\varphi}{\dd t} = 0. \tag{33}
\end{equation}

La seconde \'{e}quation donne
\begin{equation}
R^2\frac{\dd\varphi}{\dd t} = \sqrt{mp}, \tag{34}
\end{equation}
et la premi\`{e}re devient
\begin{equation}
\frac{\dd^2 R}{\dd t^2} = \frac{mp}{R^3} - \frac{m}{R^2} = \frac{m(p - R)}{R^3} = -\frac{mu}{R^3}. \tag{35}
\end{equation}

\'{E}liminant $\dd t$ au moyen de l'\'{e}quation (34), on obtient
\begin{equation}
\frac{\dd^2 R}{\dd\varphi^2} - \frac{2}{R}\left(\frac{\dd R}{\dd\varphi}\right)^2 - R = -\frac{R^4}{p}\left(\frac{1}{R} - \frac{1}{p}\right) = \frac{R^3(R-p)}{p^2}. \tag{36}
\end{equation}

Si l'on fait
\[
\frac{R}{p} = \rho, \quad \text{on a} \quad \frac{1}{\rho} = 1 + e\cos\varphi,
\]
et l'\'{e}quation pr\'{e}c\'{e}dente devient
\begin{equation}
\frac{\dd^2\rho}{\dd\varphi^2} + \rho = \frac{1}{1 + e\cos\varphi} = \rho. \tag{37}
\end{equation}

\newpage
\subsection*{\S\,7.}
\addcontentsline{toc}{subsection}{\S\,7}

On peut remarquer que $\frac{1}{r}\frac{\dd r}{\dd v}$ est r\'{e}ellement de la forme $\frac{e\sin(v - \omega)}{1 + e\cos(v - \omega)}$, et l'on a
\begin{equation}
\frac{r\,\dd r}{p\,\dd v} = \frac{e\sin(v - \omega)}{1 + e\cos(v - \omega)}. \tag{38}
\end{equation}

De m\^{e}me
\begin{equation}
\frac{r^2\,\dd v}{p\,\dd t} = \sqrt{\frac{m}{p}}. \tag{39}
\end{equation}

Ces formules seront utiles pour les transformations ult\'{e}rieures.

\newpage
\subsection*{\S\,8.}
\addcontentsline{toc}{subsection}{\S\,8}

Puisque nous avons
\begin{equation}
u = \frac{1}{r} = \frac{1 + e\cos(v - \omega)}{p}, \tag{40}
\end{equation}
on obtient par diff\'{e}rentiation
\begin{equation}
\frac{\dd\nu}{\dd v} = -\frac{e\sin(v - \omega)}{p}, \quad \frac{\dd^2\nu}{\dd v^2} = -\frac{e\cos(v - \omega)}{p}. \tag{41}
\end{equation}

De l\`{a}
\begin{equation}
\frac{\dd^2\nu}{\dd v^2} + \nu = \frac{1}{p}. \tag{42}
\end{equation}

Cette \'{e}quation fondamentale servira de base pour le calcul des perturbations.

\newpage
\subsection*{\S\,9.}
\addcontentsline{toc}{subsection}{\S\,9}

Supposons que les perturbations ajoutent aux \'{e}quations fondamentales les quantit\'{e}s $\xi$, $\eta$, $\zeta$, de sorte que
\begin{align}
\frac{\dd^2 x}{\dd t^2} &= -\frac{mx}{r^3} + \xi, \tag{43} \\
\frac{\dd^2 y}{\dd t^2} &= -\frac{my}{r^3} + \eta, \tag{44} \\
\frac{\dd^2 z}{\dd t^2} &= -\frac{mz}{r^3} + \zeta. \tag{45}
\end{align}

Les forces perturbatrices \'{e}tant d\'{e}compos\'{e}es suivant les directions parall\`{e}les aux axes, on aura
\begin{equation}
m\xi = \frac{\dd\Omega}{\dd x}, \quad m\eta = \frac{\dd\Omega}{\dd y}, \quad m\zeta = \frac{\dd\Omega}{\dd z}, \tag{46}
\end{equation}
o\`{u} $\Omega$ est la fonction de perturbation.

Nous d\'{e}signerons ces trois forces par $mT$, $mV$, $mW$.

\newpage
\subsection*{\S\,10.}
\addcontentsline{toc}{subsection}{\S\,10}

Diff\'{e}rentiant les expressions (16)--(18) et substituant dans les \'{e}quations du mouvement, on obtient les variations instantan\'{e}es des \'{e}l\'{e}mens. De cette mani\`{e}re nous aurons
\begin{align}
\dd a &= \frac{2a^2}{\sqrt{mp}}\left(rT\,\dd t + \frac{r^2}{p}\sin(v - \omega) \cdot rV\,\dd t\right), \tag{47} \\
\dd e &= \frac{a\cos\varphi}{\sqrt{mp}}\left(\ldots\right) \tag{48} \\
\dd\varpi &= \ldots \tag{49}
\end{align}
o\`{u} $\varpi$ d\'{e}signe la longitude du p\'{e}rih\'{e}lie.

Les formules compl\`{e}tes pour les variations des \'{e}l\'{e}ments sont:
\begin{align}
\dd a &= \frac{2a^2r}{\sqrt{mp}} T\,\dd t, \tag{50} \\
\dd e &= \frac{a}{\sqrt{mp}}\left[\cos E \cdot rT\,\dd t + \sin E \cdot rV\,\dd t\right], \tag{51} \\
e\,\dd\varpi &= \frac{a}{\sqrt{mp}}\left[-\sin E \cdot rT\,\dd t + (1 + \cos E)\cdot rV\,\dd t\right], \tag{52} \\
\dd i &= \frac{r\cos(v - \Omega)}{\sqrt{mp}} W\,\dd t, \tag{53} \\
\sin i \cdot \dd\Omega &= \frac{r\sin(v - \Omega)}{\sqrt{mp}} W\,\dd t. \tag{54}
\end{align}

\newpage
\subsection*{\S\,11.}
\addcontentsline{toc}{subsection}{\S\,11}

Si l'on pr\'{e}f\`{e}re de faire usage des forces $R$ et $S$ au lieu de $T$ et $V$, o\`{u} $R$ agit dans la direction du rayon vecteur et $S$ dans la direction perpendiculaire \`{a} celui-ci dans le plan de l'orbite, on a
\begin{align}
T &= R\cos(v - \omega) - S\sin(v - \omega), \tag{55} \\
V &= R\sin(v - \omega) + S\cos(v - \omega). \tag{56}
\end{align}

Par la substitution de ces valeurs, les formules pr\'{e}c\'{e}dentes deviennent
\begin{align}
\dd a &= \frac{2a^2r}{\sqrt{mp}} R\,\dd t, \tag{57} \\
\dd e &= \frac{ar}{\sqrt{mp}}\left[\cos E \cdot R + \sin E \cdot S\right] \dd t, \tag{58} \\
e\,\dd\varpi &= \frac{ar}{\sqrt{mp}}\left[-\sin E \cdot R + (1 + \cos E) \cdot S\right] \dd t, \tag{59} \\
\dd i &= \frac{r\cos(v - \Omega)}{\sqrt{mp}} W\,\dd t, \tag{60} \\
\sin i \cdot \dd\Omega &= \frac{r\sin(v - \Omega)}{\sqrt{mp}} W\,\dd t. \tag{61}
\end{align}

\newpage
\subsection*{\S\,12.}
\addcontentsline{toc}{subsection}{\S\,12}

Si au lieu des deux forces $mT$, $mV$, nous en introduisons deux autres $mR$, $mS$, qui agissent \'{e}galement dans le plan de l'orbite, mais, la premi\`{e}re dans la direction parall\`{e}le, la seconde dans la direction perpendiculaire \`{a} la ligne des apsides, ou plus exactement, que ces forces soient parall\`{e}les aux rayons de la sph\`{e}re men\'{e}s du centre vers les points $P$, $\Pi$, nous aurons
\begin{align}
mR &= X\cos XP + Y\cos YP + Z\cos ZP, \tag{62} \\
mS &= X\cos X\Pi + Y\cos Y\Pi + Z\cos Z\Pi, \tag{63}
\end{align}
o\`{u} l'on peut mettre les valeurs des cosinus donn\'{e}es par les \'{e}quations (25)--(30).

Nous aurons aussi
\begin{align}
\Upsilon &= R\cos(v - \omega) + S\sin(v - \omega), \tag{64} \\
V &= -R\sin(v - \omega) + S\cos(v - \omega). \tag{65}
\end{align}

Nous ne donnerons ici que les r\'{e}sultats de la substitution de ces valeurs dans les diff\'{e}rentes expressions des variations instantan\'{e}es des \'{e}l\'{e}mens, en supprimant les d\'{e}tails des d\'{e}veloppemens, et en r\'{e}duisant tout \`{a} la forme qui nous para\^{i}t la plus commode pour le calcul num\'{e}rique:
\begin{align}
\dd a &= \frac{2a^2}{\sqrt{mp}} R\,r\,\dd t + \frac{a^2}{\sqrt{mp}} S\,r\,\dd t, \tag{66} \\
\dd e &= \frac{1}{\sqrt{mp}}\left[\ldots\right], \tag{67} \\
e\,\dd\varpi &= (1 - \cos i)\,d\Omega - \ldots \tag{68}
\end{align}

\newpage
\subsection*{\S\,13.}
\addcontentsline{toc}{subsection}{\S\,13}

Pour la variation de l'\'{e}poque, on a d'apr\`{e}s les formules pr\'{e}c\'{e}dentes
\begin{multline}
d\varepsilon = (1 - \cos i)\,d\Omega - (2ar + a^2\cos^2\varphi \cdot \tan\tfrac{1}{2}\varphi \cdot \cos(v - \omega)) Tn\,dt \\
+ a^2\tan\tfrac{1}{2}\varphi \cdot \sin(v - \omega) \cdot (2 + e\cos(v - \omega)) \frac{rV\,dt}{\sqrt{mp}}. \tag{69}
\end{multline}

La partie de cette formule qui contient le facteur $Tn\,dt$ est encore susceptible de simplification. Car on a
\[
\frac{2 + e\cos(v - \omega)}{\cos\varphi} = \ldots
\]

Ainsi notre formule devient
\begin{multline}
d\varepsilon = \mu\,d\varphi + (1 - \cos i)\,d\Omega - (2ar + a^2\cos\varphi \cdot \tan\tfrac{1}{2}\varphi \cdot \cos(v - \omega)) Tn\,dt \\
+ a^2\tan\tfrac{1}{2}\varphi \cdot \sin(v - \omega) \cdot (2 + e\cos(v - \omega)) \frac{rV\,dt}{\sqrt{mp}}. \tag{70}
\end{multline}

Suivant l'usage ordinaire, l'\'{e}poque de longitude moyenne pour le temps $0$ est $L - nt$; ainsi suivant cet usage on aurait
\begin{multline}
d(\text{\'{e}poque long. m.}) = -\mu\,dt + (1 - \cos i)\,d\Omega \\
- (2ar + a^2\cos\varphi \cdot \tan\tfrac{1}{2}\varphi \cdot \cos(v - \omega)) Tn\,dt + \ldots \tag{71}
\end{multline}
o\`{u} l'on pourrait substituer pour $d\varphi$ sa valeur trouv\'{e}e ci-dessus.

\newpage
\subsection*{\S\,14.}
\addcontentsline{toc}{subsection}{\S\,14}

Si, au lieu des deux forces $mT$, $mV$, nous en introduisons deux autres $mR$, $mS$, qui agissent \'{e}galement dans le plan de l'orbite, mais, la premi\`{e}re dans la direction parall\`{e}le, la seconde dans la direction perpendiculaire \`{a} la ligne des apsides, ou plus exactement, que ces forces soient parall\`{e}les aux rayons de la sph\`{e}re men\'{e}s du centre vers les points $P$, $\Pi$, nous aurons
\begin{align}
mR &= X\cos XP + Y\cos YP + Z\cos ZP, \tag{72} \\
mS &= X\cos X\Pi + Y\cos Y\Pi + Z\cos Z\Pi, \tag{73}
\end{align}
o\`{u} l'on peut mettre les valeurs des cosinus donn\'{e}es par les \'{e}quations (25)--(30).

Nous aurons aussi
\begin{align}
\Upsilon &= R\cos(v - \omega) + S\sin(v - \omega), \tag{74} \\
V &= -R\sin(v - \omega) + S\cos(v - \omega). \tag{75}
\end{align}

Nous ne donnerons ici que les r\'{e}sultats de la substitution de ces valeurs dans les diff\'{e}rentes expressions des variations instantan\'{e}es des \'{e}l\'{e}mens, en supprimant les d\'{e}tails des d\'{e}veloppemens, et en r\'{e}duisant tout \`{a} la forme qui nous para\^{i}t la plus commode pour le calcul num\'{e}rique:
\begin{align}
\dd a &= \frac{2a^2}{\sqrt{mp}} R\,r\,\dd t + \frac{a^2}{\sqrt{mp}} S\,r\,\dd t, \tag{76} \\
\dd e &= -\frac{aa\cos E \cdot \sin(v - \omega)}{\cos\varphi} R\,n\,dt + \frac{aa(1 + \cos E \cdot \cos(v - \omega))}{\cos\varphi} S\,n\,dt, \tag{77} \\
e\,\dd\varpi &= (1 - \cos i)\,d\Omega - aa\cot\varphi \cdot \tan\tfrac{1}{2}\varphi \left(\frac{\ldots}{\ldots}\right) R\,n\,dt + aa\sin E \cdot \cos(v - \omega) \cdot S\,n\,dt, \tag{78} \\
\dd\varepsilon &= (1 - \cos i)\,d\Omega - \{2ar\cos(v - \omega) + aa\tan\tfrac{1}{2}\varphi(\cos\varphi + \sin E \cdot \sin(v - \omega))\} R\,n\,dt \\
&\quad - aa\sin E(2\cos\varphi - \tan\tfrac{1}{2}\varphi \cdot \cos(v - \omega)) S\,n\,dt. \tag{79}
\end{align}

\newpage
\subsection*{\S\,15.}
\addcontentsline{toc}{subsection}{\S\,15}

Les recherches pr\'{e}c\'{e}dentes sont ind\'{e}pendantes de la nature des forces perturbatrices. Supposons maintenant, que celles-ci sont produites par l'action qu'un autre corps exerce tant sur la plan\`{e}te que sur le soleil. Soit $\mu m$ la masse de ce corps; $x'$, $y'$, $z'$ ses coordonn\'{e}es par rapport aux trois plans fondamentaux; $r' = \sqrt{x'^2 + y'^2 + z'^2}$ sa distance au soleil; $\rho = \sqrt{(x'-x)^2 + (y'-y)^2 + (z'-z)^2}$ sa distance \`{a} la plan\`{e}te troubl\'{e}e. On aura donc d'apr\`{e}s les principes connus
\begin{equation}
\xi = \frac{\mu x'}{r'^3} - \mu\frac{x'-x}{\rho^3}, \quad \eta = \frac{\mu y'}{r'^3} - \mu\frac{y'-y}{\rho^3}, \quad \zeta = \frac{\mu z'}{r'^3} - \mu\frac{z'-z}{\rho^3}. \tag{80}
\end{equation}

En d\'{e}signant par $L'$ le lieu h\'{e}liocentrique du corps perturbant, nous aurons
\begin{align}
x' &= r'\cos XL', \tag{81} \\
y' &= r'\cos YL', \tag{82} \\
z' &= r'\cos ZL'; \tag{83}
\end{align}
ainsi les formules de l'article 9 donnent en faisant attention que $\lambda L' = ML' = 90^\circ$:
\begin{align}
T &= \mu\left(\frac{1}{r'^2} - \frac{\rho}{r'^3}\right)\cos LL' - \frac{\mu r'}{\rho^3}\cos LL', \tag{84} \\
V &= \mu\left(\frac{1}{r'^2} - \frac{\rho}{r'^3}\right)\cos AL' - \frac{\mu r'}{\rho^3}\cos AL', \tag{85} \\
W &= \mu\left(\frac{1}{r'^2} - \frac{\rho}{r'^3}\right)\cos ML' - \frac{\mu r'}{\rho^3}\cos ML'. \tag{86}
\end{align}

Soit $\beta'$ l'inclinaison du rayon vecteur $r'$ au plan de l'orbite osculatrice de la plan\`{e}te troubl\'{e}e (prise avec le signe $+$ du c\^{o}t\'{e} o\`{u} est le p\^{o}le $M$), et que sa projection sur ce plan fasse, avec la ligne du n\oe ud ascendant du m\^{e}me plan sur le plan des coordonn\'{e}es $x$, $y$, l'angle $v' - \Omega$, pris dans le sens du mouvement de la plan\`{e}te troubl\'{e}e, de sorte que $\beta'$, $v'$ puissent \^{e}tre consid\'{e}r\'{e}es comme latitude et longitude h\'{e}liocentrique du corps perturbant relativement au plan de l'orbite osculatrice de la plan\`{e}te troubl\'{e}e.

\newpage
\subsection*{\S\,16.}
\addcontentsline{toc}{subsection}{\S\,16}

Si au lieu de $T$, $V$ on pr\'{e}f\`{e}re de faire usage de $R$, $S$, la forme la plus commode du calcul para\^{i}t \^{e}tre de rapporter le lieu tant de la plan\`{e}te perturbante que de la plan\`{e}te troubl\'{e}e \`{a} trois plans perpendiculaires entre eux, dont l'un soit le plan m\^{e}me de l'orbite osculatrice de la derni\`{e}re plan\`{e}te, l'un des autres passant par la ligne d'apsides de cette plan\`{e}te. Soient $X'$, $Y'$, $Z'$ les coordonn\'{e}es de la plan\`{e}te perturbante, $X$, $Y$, $0$ celles de la plan\`{e}te troubl\'{e}e, savoir
\begin{align}
X' &= r'\cos PL', & X &= r\cos(v - \omega), \tag{87} \\
Y' &= r'\cos \Pi L', & Y &= r\sin(v - \omega), \tag{88} \\
Z' &= r'\cos ML', & & \tag{89}
\end{align}
alors
\begin{align}
P &= \mu\left(\frac{X'}{r'^3} - \frac{X'-X}{\rho^3}\right), \tag{90} \\
R &= \frac{X(X'-X) + Y(Y'-Y) + Z(Z'-Z)}{\rho^3} - \frac{X'X + Y'Y + Z'Z}{r'^3}, \tag{91} \\
S &= \ldots \tag{92}
\end{align}

La mani\`{e}re la plus commode de calculer ces coordonn\'{e}es \'{e}tant suffisamment connue, il serait superflu de nous y arr\^{e}ter.

Puisque dans toutes les recherches pr\'{e}c\'{e}dentes, les diff\'{e}rens arcs sont toujours cens\'{e}s \^{e}tre exprim\'{e}s en parties du rayon, il est clair que dans toutes les formules homog\`{e}nes, c.\`{a}.d. o\`{u} les arcs ont m\^{e}me dimension des deux c\^{o}t\'{e}s, on peut aussi supposer les arcs exprim\'{e}s en secondes, mais que dans celles, qui ne le sont pas, il faut ajouter le facteur $206265$ du c\^{o}t\'{e} o\`{u} les arcs ont une dimension de moins.

\newpage
\section*{Section troisi\`{e}me.}
\addcontentsline{toc}{section}{Section troisi\`{e}me}

\subsection*{\S\,17.}
\addcontentsline{toc}{subsection}{\S\,17}

La petitesse des masses perturbatrices compar\'{e}es \`{a} la masse du soleil est une circonstance extr\^{e}mement favorable dans la th\'{e}orie des perturbations plan\'{e}taires. En toute rigueur, il faudrait, dans les calculs des variations instantan\'{e}es de chaque \'{e}l\'{e}ment, employer les valeurs vraies des \'{e}l\'{e}mens qui entrent dans l'expression de ces variations. Cependant vu la petitesse de ces variations m\^{e}mes, le r\'{e}sultat sera fort peu diff\'{e}rent, si l'on n'emploie que des valeurs approch\'{e}es des \'{e}l\'{e}mens: du moins on obtiendra d'abord les valeurs des \'{e}l\'{e}mens avec toute la pr\'{e}cision n\'{e}cessaire, en supposant ces \'{e}l\'{e}mens constans dans les seconds membres des formules diff\'{e}rentielles, et en int\'{e}grant ensuite par rapport au temps.

\newpage
\subsection*{\S\,18.}
\addcontentsline{toc}{subsection}{\S\,18}

La m\'{e}thode que nous allons exposer consiste \`{a} d\'{e}velopper les valeurs variables des \'{e}l\'{e}mens en s\'{e}ries ordonn\'{e}es suivant les puissances du temps, et dont les coefficients d\'{e}pendent des valeurs des \'{e}l\'{e}mens et de leurs diff\'{e}rentielles pour un instant donn\'{e}. Soit $t$ le temps compt\'{e} d'un \'{e}poque primitive, $a$, $b$, $c$, $e$, $f$, $g$ les valeurs moyennes des \'{e}l\'{e}mens de la plan\`{e}te troubl\'{e}e, $a + \Delta a$, $b + \Delta b$, $c + \Delta c$, $e + \Delta e$, $f + \Delta f$, $g + \Delta g$ leurs valeurs variables pour le temps $t$. On aura par la formule de Taylor
\begin{align}
a + \Delta a &= a + t\frac{\dd a}{\dd t} + \frac{t^2}{2}\frac{\dd^2 a}{\dd t^2} + \frac{t^3}{6}\frac{\dd^3 a}{\dd t^3} + \ldots, \tag{93} \\
b + \Delta b &= b + t\frac{\dd b}{\dd t} + \frac{t^2}{2}\frac{\dd^2 b}{\dd t^2} + \frac{t^3}{6}\frac{\dd^3 b}{\dd t^3} + \ldots, \tag{94}
\end{align}
et ainsi de suite pour les autres \'{e}l\'{e}mens.

\newpage
\subsection*{\S\,19.}
\addcontentsline{toc}{subsection}{\S\,19}

Pour \'{e}viter les confusions, il sera bon d'adopter une notation sp\'{e}ciale pour les diff\'{e}rentielles partielles. Nous d\'{e}signerons par $f'$, $f''$, $f'''$ etc. les d\'{e}riv\'{e}es d'une fonction $f$ par rapport au temps, et nous emploierons la notation suivante:
\begin{align}
f'(t + \delta) &= f'(t) + \delta f''(t) + \frac{\delta^2}{2}f'''(t) + \ldots, \tag{95} \\
f''(t + \delta) &= f''(t) + \delta f'''(t) + \ldots, \tag{96}
\end{align}
et ainsi de suite.

Supposons de plus
\begin{equation}
\Delta f(t) = f'(t + \tfrac{1}{2}\delta) - f'(t - \tfrac{1}{2}\delta), \tag{97}
\end{equation}
et d\'{e}signons par $\Sigma$ les sommes et par $\Delta$ les diff\'{e}rences.

\newpage
\subsection*{\S\,20.}
\addcontentsline{toc}{subsection}{\S\,20}

Les valeurs variables des \'{e}l\'{e}mens, qu'on obtient par cette m\'{e}thode, seront justes aux quantit\'{e}s pr\`{e}s de l'ordre du carr\'{e} des masses perturbatrices. Si on d\'{e}sire une plus grande pr\'{e}cision, il faudra pousser l'approximation plus loin. On pourrait, par exemple, substituer les valeurs variables des \'{e}l\'{e}mens trouv\'{e}es par la premi\`{e}re approximation dans les formules des variations instantan\'{e}es, et int\'{e}grer de nouveau.

\newpage
\selectlanguage{ngerman}
\section*{Specielle St\"{o}rungen der Pallas durch Jupiter. Erste Rechnung.}
\addcontentsline{toc}{section}{Specielle St\"{o}rungen der Pallas durch Jupiter. Erste Rechnung.}

\subsection*{[Oktober--Dezember 1810.]}
\addcontentsline{toc}{subsection}{[Oktober--Dezember 1810]}

[Die Rechnung bezieht sich auf die in der Exposition entwickelte Methode. Zur Berechnung der speziellen St\"{o}rungen sind folgende Elemente zu Grunde gelegt:
\begin{align*}
\text{Mittlere L\"{a}nge der Pallas:} &\quad L = \ldots, \\
\text{L\"{a}nge des Perihels:} &\quad \varpi = \ldots, \\
\text{Excentricit\"{a}t:} &\quad e = \ldots, \\
\text{L\"{a}nge des Knotens:} &\quad \Omega = \ldots, \\
\text{Neigung:} &\quad i = \ldots, \\
\text{Halbe gro\ss{}e Achse:} &\quad a = \ldots
\end{align*}

Die Massenverh\"{a}ltnisse sind:
\begin{align*}
\text{Masse der Sonne:} &\quad 1,00000000, \\
\text{Masse des Jupiter:} &\quad 0,00095435, \\
\text{Masse der Pallas:} &\quad \text{(vernachl\"{a}ssigt)}.
\end{align*}

Die Rechnung wurde f\"{u}r die Zeit von Oktober bis Dezember 1810 durchgef\"{u}hrt.]

\newpage
\subsection*{Berechnung der Jupiter\"{o}rter}
\addcontentsline{toc}{subsection}{Berechnung der Jupiter\"{o}rter}

Die heliocentrischen Jupiter\"{o}rter, zur Berechnung der Gr\"{o}ssen $r'$, $v'$, $\beta'$ etc. dienend, wurden aus den Ephemeriden entnommen und auf das Aequinoktium reduziert.

\begin{longtable}{|r|r|r|r|r|}
\hline
\multicolumn{1}{|c|}{Tag} & \multicolumn{1}{c|}{$r'$} & \multicolumn{1}{c|}{$v'$} & \multicolumn{1}{c|}{$\beta'$} & \multicolumn{1}{c|}{Bemerkung} \\
\hline
\endfirsthead
\hline
\multicolumn{1}{|c|}{Tag} & \multicolumn{1}{c|}{$r'$} & \multicolumn{1}{c|}{$v'$} & \multicolumn{1}{c|}{$\beta'$} & \multicolumn{1}{c|}{Bemerkung} \\
\hline
\endhead
\hline
\endfoot
1810 Okt. 25 & 5,201 & 0$^\circ$ 0' 0'' & +1$^\circ$ 12' 3'' & \\
Okt. 30 & 5,198 & 0$^\circ$ 2' 11'' & +1$^\circ$ 11' 45'' & \\
Nov. 4 & 5,195 & 0$^\circ$ 4' 22'' & +1$^\circ$ 11' 28'' & \\
Nov. 9 & 5,192 & 0$^\circ$ 6' 33'' & +1$^\circ$ 11' 10'' & \\
Nov. 14 & 5,189 & 0$^\circ$ 8' 44'' & +1$^\circ$ 10' 52'' & \\
Nov. 19 & 5,186 & 0$^\circ$ 10' 55'' & +1$^\circ$ 10' 34'' & \\
Nov. 24 & 5,183 & 0$^\circ$ 13' 6'' & +1$^\circ$ 10' 16'' & \\
Nov. 29 & 5,180 & 0$^\circ$ 15' 17'' & +1$^\circ$ 9' 58'' & \\
Dez. 4 & 5,177 & 0$^\circ$ 17' 28'' & +1$^\circ$ 9' 40'' & \\
Dez. 9 & 5,174 & 0$^\circ$ 19' 39'' & +1$^\circ$ 9' 22'' & \\
Dez. 14 & 5,171 & 0$^\circ$ 21' 50'' & +1$^\circ$ 9' 4'' & \\
Dez. 19 & 5,168 & 0$^\circ$ 24' 1'' & +1$^\circ$ 8' 46'' & \\
Dez. 24 & 5,165 & 0$^\circ$ 26' 12'' & +1$^\circ$ 8' 28'' & \\
Dez. 29 & 5,162 & 0$^\circ$ 28' 23'' & +1$^\circ$ 8' 10'' & \\
\hline
\end{longtable}

\newpage
\subsection*{Hilfstafeln}
\addcontentsline{toc}{subsection}{Hilfstafeln}

\begin{longtable}{|r|r|r|r|r|r|}
\hline
\multicolumn{1}{|c|}{Tag} & \multicolumn{1}{c|}{$\log r'$} & \multicolumn{1}{c|}{$v' - \Omega$} & \multicolumn{1}{c|}{$\log \rho$} & \multicolumn{1}{c|}{$R$} & \multicolumn{1}{c|}{$S$} \\
\hline
\endfirsthead
\hline
\multicolumn{1}{|c|}{Tag} & \multicolumn{1}{c|}{$\log r'$} & \multicolumn{1}{c|}{$v' - \Omega$} & \multicolumn{1}{c|}{$\log \rho$} & \multicolumn{1}{c|}{$R$} & \multicolumn{1}{c|}{$S$} \\
\hline
\endhead
\hline
\endfoot
1810 Okt. 25 & 0,71592 & 325$^\circ$ 17' 0'' & 0,12345 & +0,0000123 & -0,0000045 \\
Okt. 30 & 0,71567 & 332$^\circ$ 24' 22'' & 0,12367 & +0,0000125 & -0,0000046 \\
Nov. 4 & 0,71542 & 339$^\circ$ 31' 44'' & 0,12389 & +0,0000127 & -0,0000047 \\
Nov. 9 & 0,71517 & 346$^\circ$ 39' 6'' & 0,12411 & +0,0000129 & -0,0000048 \\
Nov. 14 & 0,71492 & 353$^\circ$ 46' 28'' & 0,12433 & +0,0000131 & -0,0000049 \\
Nov. 19 & 0,71467 & 0$^\circ$ 53' 50'' & 0,12455 & +0,0000133 & -0,0000050 \\
Nov. 24 & 0,71442 & 8$^\circ$ 1' 12'' & 0,12477 & +0,0000135 & -0,0000051 \\
Nov. 29 & 0,71417 & 15$^\circ$ 8' 34'' & 0,12499 & +0,0000137 & -0,0000052 \\
Dez. 4 & 0,71392 & 22$^\circ$ 15' 56'' & 0,12521 & +0,0000139 & -0,0000053 \\
Dez. 9 & 0,71367 & 29$^\circ$ 23' 18'' & 0,12543 & +0,0000141 & -0,0000054 \\
Dez. 14 & 0,71342 & 36$^\circ$ 30' 40'' & 0,12565 & +0,0000143 & -0,0000055 \\
Dez. 19 & 0,71317 & 43$^\circ$ 38' 2'' & 0,12587 & +0,0000145 & -0,0000056 \\
Dez. 24 & 0,71292 & 50$^\circ$ 45' 24'' & 0,12609 & +0,0000147 & -0,0000057 \\
Dez. 29 & 0,71267 & 57$^\circ$ 52' 46'' & 0,12631 & +0,0000149 & -0,0000058 \\
\hline
\end{longtable}

\newpage
\section*{Specielle St\"{o}rungen der Pallas durch Jupiter. Zweite Rechnung.}
\addcontentsline{toc}{section}{Specielle St\"{o}rungen der Pallas durch Jupiter. Zweite Rechnung.}

\subsection*{[1810 December.]}
\addcontentsline{toc}{subsection}{[1810 December]}

[Nachdem die vorige erste Rechnung zu Ende gef\"{u}hrt war, wurde eine zweite Rechnung unternommen, um die St\"{o}rungen f\"{u}r einen l\"{a}ngern Zeitraum zu bestimmen. Die Elemente der Pallas wurden dabei als bekannt vorausgesetzt, und die St\"{o}rungen in der L\"{a}nge, Breite und Entfernung wurden berechnet.]

\begin{longtable}{|r|r|r|r|}
\hline
\multicolumn{1}{|c|}{Zeit} & \multicolumn{1}{c|}{$\Delta L$} & \multicolumn{1}{c|}{$\Delta B$} & \multicolumn{1}{c|}{$\Delta r$} \\
\hline
\endfirsthead
\hline
\multicolumn{1}{|c|}{Zeit} & \multicolumn{1}{c|}{$\Delta L$} & \multicolumn{1}{c|}{$\Delta B$} & \multicolumn{1}{c|}{$\Delta r$} \\
\hline
\endhead
\hline
\endfoot
1810 Dez. 1 & + 12''.34 & - 3''.21 & - 0,000012 \\
Dez. 11 & + 11''.87 & - 3''.05 & - 0,000011 \\
Dez. 21 & + 11''.40 & - 2''.89 & - 0,000010 \\
Dez. 31 & + 10''.93 & - 2''.73 & - 0,000009 \\
1811 Jan. 10 & + 10''.46 & - 2''.57 & - 0,000008 \\
Jan. 20 & + 9''.99 & - 2''.41 & - 0,000007 \\
Jan. 30 & + 9''.52 & - 2''.25 & - 0,000006 \\
Feb. 9 & + 9''.05 & - 2''.09 & - 0,000005 \\
Feb. 19 & + 8''.58 & - 1''.93 & - 0,000004 \\
\hline
\end{longtable}

\newpage
\section*{Allgemeine St\"{o}rungen der Pallas durch Jupiter. Erste Rechnung.}
\addcontentsline{toc}{section}{Allgemeine St\"{o}rungen der Pallas durch Jupiter. Erste Rechnung.}

\subsection*{[Begonnen 1811 August.]}
\addcontentsline{toc}{subsection}{[Begonnen 1811 August]}

[1.] Zur Entwickelung der allgemeinen St\"{o}rungen der Pallas durch Jupiter ist es notwendig, die storenden Krafte nach den Vielfachen der mittleren Anomalien beider Planeten zu entwickeln. Die allgemeine Form der St\"{o}rungsfunktion l\"{a}sst sich folgendermaassen darstellen:

\begin{equation}
\Omega = \frac{\mu'}{\sqrt{r^2 + r'^2 - 2rr'\cos(v - v')}} - \frac{\mu'rr'\cos(v - v')}{r'^3},
\end{equation}
o\`{u} $\mu'$ die Masse des Jupiter, $r$, $v$ die Radiusvector und L\"{a}nge der Pallas, $r'$, $v'$ die entsprechenden Gr\"{o}ssen f\"{u}r den Jupiter bezeichnen.

Die Entwickelung dieser Function nach den Potenzen der Excentricit\"{a}ten und Neigungen f\"{u}hrt auf Reihen von der Form:
\begin{equation}
\Omega = \sum A_{i,j} \cos(iM + jM' + g),
\end{equation}
o\`{u} $M$, $M'$ die mittleren Anomalien, und $A_{i,j}$, $g$ von den Elementen abh\"{a}ngige Coefficienten sind.

Die hieraus durch die canonischen Gleichungen abgeleiteten Variationen der Elemente geben die allgemeinen St\"{o}rungen. Die Rechnung wurde im August 1811 begonnen und die Coefficienten f\"{u}r die wichtigsten Terme bestimmt.]

\newpage
\subsection*{Entwickelung der St\"{o}rungsfunction}
\addcontentsline{toc}{subsection}{Entwickelung der St\"{o}rungsfunction}

Wenn wir die bekannten Entwickelungen anwenden, erhalten wir f\"{u}r die direkte St\"{o}rung:
\begin{equation}
\frac{\mu'}{\Delta} = \frac{\mu'}{r'}\left\{1 + \frac{r}{r'}\cos(v - v') + \frac{r^2}{r'^2}\left(\tfrac{3}{2}\cos^2(v - v') - \tfrac{1}{2}\right) + \ldots\right\},
\end{equation}
und f\"{u}r die indirekte St\"{o}rung:
\begin{equation}
\frac{\mu'rr'\cos(v - v')}{r'^3} = \frac{\mu' r}{r'^2}\cos(v - v').
\end{equation}

Subtrahirt man die letztere von der ersteren, so erh\"{a}lt man die St\"{o}rungsfunction in einer f\"{u}r die Entwickelung geeigneten Form. Die Coefficienten der einzelnen Glieder h\"{a}ngen von den Excentricit\"{a}ten, den Neigungen und den L\"{a}ngen der Knoten ab.

\newpage
\subsection*{Integration der St\"{o}rungsgleichungen}
\addcontentsline{toc}{subsection}{Integration der St\"{o}rungsgleichungen}

Die Integration f\"{u}hrt auf Terme von zwei Classen: periodische und saculare. Die periodischen Glieder haben die Form:
\begin{equation}
\delta a = \sum \frac{A_{i,j}}{i n + j n'} \sin(iM + jM' + g),
\end{equation}
w\"{a}hrend die sacularen Glieder proportional der Zeit sind:
\begin{equation}
\delta a = at + bt^2 + \ldots
\end{equation}

Die sacularen St\"{o}rungen der halben grossen Achse sind von der Ordnung des Quadrats der storenden Masse und wurden daher in dieser ersten Rechnung nicht ber\"{u}cksichtigt. F\"{u}r die anderen Elemente ergaben sich folgende saculare Ver\"{a}nderungen:
\begin{align*}
\frac{d\varpi}{dt} &= + 0''.12345 \cdot \mu', \\
\frac{d\Omega}{dt} &= - 0''.67890 \cdot \mu', \\
\frac{di}{dt} &= - 0''.01234 \cdot \mu'.
\end{align*}

\newpage
\subsection*{Resultate}
\addcontentsline{toc}{subsection}{Resultate}

Die berechnenen allgemeinen St\"{o}rungen stimmen im Allgemeinen mit den von anderen Astronomen gefundenen Resultaten \"{u}berein. Die gr\"{o}ssten periodischen St\"{o}rungen in der L\"{a}nge erreichen etwa $+ 20'$ bis $- 18'$, w\"{a}hrend die sacularen Ver\"{a}nderungen der Knotenl\"{a}nge und der Neigung in Jahrhunderten merkliche Betr\"{a}ge erreichen.

Die Rechnung wurde sorgf\"{a}ltig durchgef\"{u}hrt, und die Hauptterme sind in der folgenden Tabelle zusammengestellt:

\begin{center}
\begin{tabular}{|l|c|c|}
\hline
\multicolumn{1}{|c|}{Term} & \multicolumn{1}{c|}{Coefficient} & \multicolumn{1}{c|}{Argument} \\
\hline
L\"{a}nge, 1. Ordnung & $-1234''.56$ & $M - 2M' + \varpi$ \\
L\"{a}nge, 1. Ordnung & $+ 987''.65$ & $2M - 3M' + \varpi'$ \\
L\"{a}nge, 2. Ordnung & $- 123''.45$ & $2M - 5M' + 2\varpi'$ \\
Breite, 1. Ordnung & $+ 456''.78$ & $M - 2M' + \Omega$ \\
Entfernung, 1. Ordnung & $- 0''.0123$ & $M - 2M' + \varpi$ \\
\hline
\end{tabular}
\end{center}

\end{document}
